\documentclass[12pt]{article}
\usepackage[utf8]{inputenc}
\usepackage[T1]{fontenc}
\usepackage[margin=1in]{geometry}
\usepackage{setspace}
\usepackage{xcolor}
\usepackage[colorlinks=true, urlcolor=blue!60!black, linkcolor=blue!60!black]{hyperref}
\setlength{\parindent}{2em}
\setlength{\parskip}{1.25em}
\renewcommand{\baselinestretch}{1.0}
\begin{document}
\begin{center}
{\LARGE \textit{Loving v. Virginia}, 388 U.S. 1 (1967)\par}
\vspace{0.8em}
{\large Case Brief\par}
\end{center}
\vspace{1em}

\section{Memory Jogger}
A unanimous Supreme Court struck down Virginia's ban on interracial marriage, holding that racial restrictions on marriage violate both the Equal Protection Clause and the substantive liberty protected by the Due Process Clause of the Fourteenth Amendment (\href{https://www.courtlistener.com/opinion/107480/loving-v-virginia/}{CourtListener}; \href{https://supreme.justia.com/cases/federal/us/388/1/}{Justia}).

\section{Detailed Case Facts}
In June 1958, Mildred Jeter, a Black woman, and Richard Loving, a white man, both residents of Caroline County, Virginia, traveled to Washington, D.C., where they were lawfully married on June 2, 1958 (\href{https://supreme.justia.com/cases/federal/us/388/1/}{388 U.S. at 2}; \href{https://encyclopediavirginia.org/entries/loving-v-virginia-1967/}{Encyclopedia Virginia}). Shortly thereafter, the couple returned to Caroline County to establish their marital home. In July 1958, local police raided their residence at night and arrested them after Mildred pointed to their District of Columbia marriage certificate on the bedroom wall (\href{https://supreme.justia.com/cases/federal/us/388/1/}{Justia}).

At the October Term, 1958, a Caroline County Circuit Court grand jury indicted the Lovings for violating Virginia's comprehensive anti-miscegenation statutory scheme, which dated to the \href{https://tile.loc.gov/storage-services/service/ll/usrep/usrep388/usrep388001/usrep388001.pdf}{Racial Integrity Act of 1924} (\href{https://supreme.justia.com/cases/federal/us/388/1/}{388 U.S. at 2--6}). Specifically, they were charged under Virginia Code \S~20-58 (prohibiting any ``white person and colored person'' from leaving Virginia to evade its marriage ban and returning to cohabit as husband and wife) and \S~20-59 (making interracial marriage a felony punishable by one to five years in the penitentiary). Related provisions automatically voided interracial marriages without judicial decree (\S~20-57) and defined ``white person'' as having no ascertainable non-Caucasian blood---with a narrow exception for persons with one-sixteenth or less American Indian blood, intended to accommodate descendants of Pocahontas and John Rolfe---while defining ``colored person'' as anyone with any ascertainable ``Negro blood'' (\S\S~20-54, 1-14) (\href{https://tile.loc.gov/storage-services/service/ll/usrep/usrep388/usrep388001/usrep388001.pdf}{388 U.S. at 4--5 \& nn.3--4}). At the time of the decision, Virginia was one of sixteen states that still prohibited and punished interracial marriage (\href{https://supreme.justia.com/cases/federal/us/388/1/}{388 U.S. at 6 \& n.5}).

\section{Procedural History}
On January 6, 1959, the Lovings pleaded guilty in the Circuit Court of Caroline County. Trial Judge Leon M. Bazile sentenced each of them to one year in jail, but suspended the sentence for twenty-five years on the condition that the couple leave Virginia and not return together for twenty-five years, writing that ``Almighty God created the races white, black, yellow, malay and red, and he placed them on separate continents,'' and that ``[t]he fact that he separated the races shows that he did not intend for the races to mix'' (\href{https://supreme.justia.com/cases/federal/us/388/1/}{388 U.S. at 3}).

After relocating to the District of Columbia, the Lovings filed a motion in the Caroline County Circuit Court on November 6, 1963, to vacate the judgment and set aside their sentences as unconstitutional under the Fourteenth Amendment (\href{https://supreme.justia.com/cases/federal/us/388/1/}{388 U.S. at 3}). When the state trial court had still not ruled by October 28, 1964, they filed a federal class action in the United States District Court for the Eastern District of Virginia seeking a three-judge court to invalidate the statutes and enjoin enforcement. On January 22, 1965, Judge Bazile denied the state motion to vacate, and the Lovings appealed to the Supreme Court of Appeals of Virginia; on February 11, 1965, the three-judge federal district court continued the federal action pending state appellate review (\href{https://supreme.justia.com/cases/federal/us/388/1/}{388 U.S. at 3}).

On March 7, 1966, the Supreme Court of Appeals of Virginia upheld the constitutionality of the anti-miscegenation statutes in \textit{Loving v. Commonwealth}, 206 Va. 924, 147 S.E.2d 78 (1966), relying on its prior decision in \textit{Naim v. Naim}, 197 Va. 80, 87 S.E.2d 749 (1955), though it modified the sentence by striking the banishment condition as unreasonable and remanding for resentencing while affirming the convictions (\href{https://supreme.justia.com/cases/federal/us/388/1/}{388 U.S. at 3--4, 7}; \href{https://co.caroline.va.us/308/The-Lovings}{Caroline County VA}). The Lovings appealed to the U.S. Supreme Court, which noted probable jurisdiction on December 12, 1966 (385 U.S. 986), heard oral argument on April 10, 1967, and issued its decision on June 12, 1967 (\href{https://www.oyez.org/cases/1966/395}{Oyez}; \href{https://supreme.justia.com/cases/federal/us/388/1/}{388 U.S. at 1, 4}).

\section{Judicial Votes}
\textbf{Vote: 9--0} (reversing the Supreme Court of Appeals of Virginia) (\href{https://www.oyez.org/cases/1966/395}{Oyez}; \href{https://supreme.justia.com/cases/federal/us/388/1/}{Justia}).
\begin{itemize}
\item \textbf{Majority Opinion (8 Justices):} Chief Justice Earl Warren delivered the opinion of the Court, joined by Justices Hugo L. Black, William O. Douglas, Tom C. Clark, John Marshall Harlan II, William J. Brennan, Jr., Byron R. White, and Abe Fortas (\href{https://supreme.justia.com/cases/federal/us/388/1/}{Justia}).
\item \textbf{Concurring in the Judgment (1 Justice):} Justice Potter Stewart filed a brief separate opinion concurring in the judgment (\href{https://supreme.justia.com/cases/federal/us/388/1/}{388 U.S. at 13}).
\item \textbf{Dissents:} None.
\end{itemize}

\section{Holding}
\textbf{Judgment:} Reversed; the Lovings' convictions under Virginia's anti-miscegenation statutes were set aside (\href{https://supreme.justia.com/cases/federal/us/388/1/}{388 U.S. at 12}).

\textbf{Rule Announced:} State statutory schemes that prohibit or punish marriage on the basis of racial classifications violate both (1) the Equal Protection Clause of the Fourteenth Amendment, because racial classifications---even if ``equally applied'' to participants of each race---are subject to ``the most rigid scrutiny'' and cannot stand absent a permissible, overriding state objective independent of invidious racial discrimination, and (2) the Due Process Clause of the Fourteenth Amendment, because the freedom to marry is a ``vital personal right'' and fundamental liberty that cannot be denied on the unsupportable basis of racial classifications (\href{https://supreme.justia.com/cases/federal/us/388/1/}{388 U.S. at 11--12}).

\textbf{Holding vs. Dicta:} The Court's twin holdings are that Virginia's racial ban on marriage violates both the Equal Protection Clause and the Due Process Clause of the Fourteenth Amendment, and that ``equal application'' does not exempt a racial classification from strict scrutiny (explicitly repudiating \textit{Pace v. Alabama}, 106 U.S. 583 (1883)). The Court's broader characterization of marriage as ``one of the vital personal rights essential to the orderly pursuit of happiness by free men'' and ``one of the `basic civil rights of man,' fundamental to our very existence and survival'' (\href{https://supreme.justia.com/cases/federal/us/388/1/}{388 U.S. at 12}) sits at the boundary between holding and dicta: while necessary to establish marriage as a fundamental liberty for substantive due process purposes, the Court tied its Due Process holding closely to the invidious racial basis of the restriction, leaving open for later cases (such as \textit{Zablocki v. Redhail}, \textit{Turner v. Safley}, and \textit{Obergefell v. Hodges}) how broadly that fundamental-right principle applies to non-racial marriage restrictions (\href{https://civics.supremecourthistory.org/article/loving-v-virginia/}{Supreme Court Historical Society}). Additionally, the Court expressly declined to rest its holding on the asymmetry of Virginia's statute (which banned only marriages involving white persons while permitting non-white groups to intermarry), noting in Footnote 11 that racial bans on marriage would violate the Fourteenth Amendment ``even assuming an even-handed state purpose to protect the `integrity' of all races'' (\href{https://supreme.justia.com/cases/federal/us/388/1/}{388 U.S. at 11 n.11}).

\section{Analysis of Opinions}

\subsection{Majority Opinion (Chief Justice Warren)}
Chief Justice Warren structured the Court's opinion into two independent constitutional grounds: Equal Protection (Part I) and Substantive Due Process (Part II) (\href{https://supreme.justia.com/cases/federal/us/388/1/}{388 U.S. at 7--12}).

In Part I, the Court first acknowledged that marriage is traditionally regulated under state police power (\textit{Maynard v. Hill}, 125 U.S. 190 (1888)), but rejected the Virginia court's Tenth Amendment theory because state marriage regulations remain bounded by the Fourteenth Amendment (\textit{Meyer v. Nebraska}, 262 U.S. 390 (1923); \textit{Skinner v. Oklahoma}, 316 U.S. 535 (1942)). Virginia principally argued an ``equal application'' theory drawn from \textit{Pace v. Alabama}, 106 U.S. 583 (1883): because Virginia punished white and Black spouses identically, the law did not invidiously discriminate and should receive only deferential rational-basis review, under which conflicting scientific evidence favored legislative deference (\href{https://supreme.justia.com/cases/federal/us/388/1/}{388 U.S. at 7--8}). The Court flatly rejected this argument. Following \textit{McLaughlin v. Florida}, 379 U.S. 184 (1964), Warren explained that ``the fact of equal application does not immunize the statute from the very heavy burden of justification which the Fourteenth Amendment has traditionally required of state statutes drawn according to race'' (\href{https://supreme.justia.com/cases/federal/us/388/1/}{388 U.S. at 9}). Addressing Virginia's reliance on debates in the Thirty-ninth Congress over the Freedmen's Bureau Bill and the Civil Rights Act of 1866, Warren invoked \textit{Brown v. Board of Education}, 347 U.S. 483, 489 (1954), concluding that those historical sources were ``at best, inconclusive'' and could not override the Fourteenth Amendment's ``clear and central purpose'' to eliminate all official state sources of invidious racial discrimination (\href{https://supreme.justia.com/cases/federal/us/388/1/}{388 U.S. at 9--10}). Applying the ``most rigid scrutiny'' (\textit{Korematsu v. United States}, 323 U.S. 214, 216 (1944)), the Court found ``patently no legitimate overriding purpose independent of invidious racial discrimination,'' pointing out that Virginia prohibited only interracial marriages involving white persons---proving the law was a measure ``designed to maintain White Supremacy'' (\href{https://supreme.justia.com/cases/federal/us/388/1/}{388 U.S. at 11--12}).

In Part II, in two concise but landmark paragraphs, the Court held that the statutes independently violated the Due Process Clause of the Fourteenth Amendment. Citing \textit{Skinner v. Oklahoma}, 316 U.S. 535, 541 (1942), Warren declared that ``[t]he freedom to marry has long been recognized as one of the vital personal rights essential to the orderly pursuit of happiness by free men'' and ``one of the `basic civil rights of man,' fundamental to our very existence and survival'' (\href{https://supreme.justia.com/cases/federal/us/388/1/}{388 U.S. at 12}). Denying that fundamental liberty on the unsupportable basis of racial classifications ``is surely to deprive all the State's citizens of liberty without due process of law.'' Thus, ``the freedom to marry, or not marry, a person of another race resides with the individual, and cannot be infringed by the State'' (\href{https://supreme.justia.com/cases/federal/us/388/1/}{388 U.S. at 12}).

\subsection{Concurring Opinion (Justice Stewart)}
Justice Stewart wrote a two-sentence concurrence in the judgment, adhering to the bright-line rule he had articulated in \textit{McLaughlin v. Florida}, 379 U.S. 184, 198 (1964): ``it is simply not possible for a state law to be valid under our Constitution which makes the criminality of an act depend upon the race of the actor'' (\href{https://supreme.justia.com/cases/federal/us/388/1/}{388 U.S. at 13}). Rather than applying strict scrutiny's means-ends balancing or invoking substantive due process, Stewart treated race-based criminal liability as per se unconstitutional.

\subsection{Dissenting Opinions}
None. The decision was unanimous (\href{https://www.oyez.org/cases/1966/395}{Oyez}).

\section{Examples: Future Applications}
\begin{enumerate}
\item \textbf{Same-Side (Unconstitutional):} A state enacts a statute prohibiting individuals of different national origins or ethnicities from marrying, asserting that the law applies ``equally'' to all groups to preserve cultural heritage. \textit{Why:} Under \textit{Loving}, ancestry and race classifications in marriage are subject to the most rigid scrutiny regardless of equal application, and preserving racial or ethnic ``purity'' is not a permissible state objective independent of invidious discrimination (\href{https://supreme.justia.com/cases/federal/us/388/1/}{388 U.S. at 11 \& n.11}).
\item \textbf{Same-Side (Unconstitutional):} A state bars incarcerated persons from marrying unless the prison superintendent finds ``compelling reasons'' (such as pregnancy) to approve the marriage, without tying the restriction to legitimate penological security needs. \textit{Why:} Building directly on \textit{Loving}'s Part II recognition of marriage as a fundamental liberty under the Due Process Clause, an arbitrary or overbroad restriction on the personal choice to marry deprives citizens of substantive liberty without sufficient justification (as later reflected in \textit{Turner v. Safley}, 482 U.S. 78 (1987)).
\item \textbf{Opposite-Side (Constitutional):} A state requires prospective spouses to obtain a marriage license, meet a minimum age of eighteen (or obtain parental/judicial consent at sixteen), and refrain from marrying close blood relatives such as siblings. \textit{Why:} \textit{Loving} reaffirmed \textit{Maynard v. Hill} that marriage remains a social relation subject to reasonable state police-power regulation (\href{https://supreme.justia.com/cases/federal/us/388/1/}{388 U.S. at 7}); non-invidious eligibility rules tied to capacity, consent, and public health do not rest on suspect classifications or arbitrarily deny the fundamental right to marry.
\item \textbf{Opposite-Side (Constitutional):} A state imposes a modest, generally applicable \$50 administrative fee and a 72-hour waiting period for issuing a marriage license, with a fee waiver for indigent applicants. \textit{Why:} Unlike Virginia's absolute criminal prohibition, a minor procedural regulation that neither discriminates on a suspect basis nor substantially interferes with the decision to marry does not trigger strict scrutiny under \textit{Loving}'s Equal Protection or Due Process holdings.
\item \textbf{Fence-Sitter:} A state prohibits any person who owes court-ordered child support from obtaining a marriage license without proving that the support obligation is current and that the children are unlikely to become public charges. \textit{Why:} Unlike \textit{Loving}, this law involves no racial classification and serves a legitimate, non-invidious state interest in child welfare. Yet under \textit{Loving}'s Part II Due Process holding, it directly and substantially burdens the fundamental right to marry for indigent noncustodial parents, creating a close doctrinal tension between state domestic-relations authority and heightened scrutiny for direct burdens on marriage (ultimately resolved against the state in \textit{Zablocki v. Redhail}, 434 U.S. 374 (1978)).
\end{enumerate}

\section{Critique}
\textbf{Originalist and Textualist Perspectives:} \textit{Loving} has long served as a central litmus test for originalism. Critics of originalism---including Richard Posner, Cass Sunstein, Jack Balkin, and Michael Klarman---have argued that originalism struggles to justify \textit{Loving} because numerous members of the Thirty-ninth Congress and state ratifiers in 1868 stated that the Fourteenth Amendment would not invalidate state anti-miscegenation laws (\href{https://scholarlycommons.law.northwestern.edu/cgi/viewcontent.cgi?article=1205&context=facultyworkingpapers}{Calabresi \& Matthews, Northwestern Working Paper}). In response, prominent originalists such as Steven G. Calabresi and Andrea Matthews have defended \textit{Loving} on original public meaning grounds, arguing that the text of the Fourteenth Amendment's Equal Protection and Privileges or Immunities Clauses constitutionalized the Civil Rights Act of 1866's guarantee of equal civil rights---including the right to make and enforce contracts, which encompassed marriage---and prohibited all racial caste legislation regardless of individual framers' inconsistent expectations (\href{https://digitalcommons.law.byu.edu/lawreview/vol2012/iss5/1/}{Calabresi \& Matthews, 2012 BYU L. Rev. 1393}; see also \href{https://virginialawreview.org/authors/ilan-wurman/}{Ilan Wurman, Va. L. Rev.}). Conservative and textualist jurists, however, frequently distinguish \textit{Loving}'s unassailable Equal Protection holding from its brief Part II Substantive Due Process passage; for instance, commentators note that while even strict originalists accept \textit{Loving}'s invalidation of racial marriage bans under the Equal Protection Clause, Justice Clarence Thomas's skepticism of substantive due process reflects an ongoing textualist critique of deriving unenumerated rights from the Due Process Clause (\href{https://reason.com/volokh/2022/06/28/justice-thomas-and-loving-v-virginia/}{Reason / Volokh Conspiracy}).

\textbf{Progressive and Doctrinal Perspectives:} Progressive and doctrinal scholars praise \textit{Loving} both for dismantling the final pillar of de jure Jim Crow racial caste and for articulating a synergistic vision of liberty and equality (\href{https://digitalcommons.nyls.edu/cgi/viewcontent.cgi?article=1091&context=nyls_law_review}{NYLS Law Review}; \href{https://scholarship.law.edu/cgi/viewcontent.cgi?article=3486&context=lawreview}{Catholic U. L. Rev.}). At the same time, some progressive scholars observe that \textit{Loving}'s Part II substantive due process analysis is remarkably brief and under-theorized---leaving later courts to debate whether the fundamental right to marry was strictly parasitic on racial equality or an autonomous right of personal choice that extends to same-sex couples, as the Court ultimately held in \textit{Obergefell v. Hodges}, 576 U.S. 644 (2015) (\href{https://civics.supremecourthistory.org/article/loving-v-virginia/}{Supreme Court Historical Society}; \href{https://scholar.smu.edu/cgi/viewcontent.cgi?article=4964&context=smulr}{SMU Law Review}). Critical race and civil rights scholars also note that emphasizing \textit{Loving} solely as an individual-autonomy case can obscure its core anti-subordination holding: the Court explicitly named and condemned ``White Supremacy'' as an unconstitutional state objective (\href{https://digitalcommons.nyls.edu/cgi/viewcontent.cgi?article=1091&context=nyls_law_review}{NYLS Law Review}).

\textbf{Underlying Competing Values:} The case squarely pits federalism and traditional state police-power control over domestic relations against individual autonomy, intimate association, and the Fourteenth Amendment's anti-caste mandate of racial equality. It also highlights the enduring methodological tension between original expected applications and the broader principles encoded in constitutional text.

\section{Key Quotations}
\begin{quote}
``At the very least, the Equal Protection Clause demands that racial classifications, especially suspect in criminal statutes, be subjected to the `most rigid scrutiny,' \textit{Korematsu v. United States}, 323 U.S. 214, 216 (1944), and, if they are ever to be upheld, they must be shown to be necessary to the accomplishment of some permissible state objective, independent of the racial discrimination which it was the object of the Fourteenth Amendment to eliminate.''
\end{quote}
\noindent --- Chief Justice Warren, Opinion of the Court, 388 U.S. at 11 (\href{https://supreme.justia.com/cases/federal/us/388/1/}{Justia}; \href{https://www.courtlistener.com/opinion/107480/loving-v-virginia/}{CourtListener}).

\begin{quote}
``There is patently no legitimate overriding purpose independent of invidious racial discrimination which justifies this classification. The fact that Virginia prohibits only interracial marriages involving white persons demonstrates that the racial classifications must stand on their own justification, as measures designed to maintain White Supremacy.''
\end{quote}
\noindent --- Chief Justice Warren, Opinion of the Court, 388 U.S. at 11 (\href{https://supreme.justia.com/cases/federal/us/388/1/}{Justia}; \href{https://www.courtlistener.com/opinion/107480/loving-v-virginia/}{CourtListener}).

\begin{quote}
``The freedom to marry has long been recognized as one of the vital personal rights essential to the orderly pursuit of happiness by free men.''
\end{quote}
\noindent --- Chief Justice Warren, Opinion of the Court, 388 U.S. at 12 (\href{https://supreme.justia.com/cases/federal/us/388/1/}{Justia}; \href{https://www.courtlistener.com/opinion/107480/loving-v-virginia/}{CourtListener}).

\begin{quote}
``Under our Constitution, the freedom to marry, or not marry, a person of another race resides with the individual, and cannot be infringed by the State.''
\end{quote}
\noindent --- Chief Justice Warren, Opinion of the Court, 388 U.S. at 12 (\href{https://supreme.justia.com/cases/federal/us/388/1/}{Justia}; \href{https://www.courtlistener.com/opinion/107480/loving-v-virginia/}{CourtListener}).

\begin{quote}
``I have previously expressed the belief that `it is simply not possible for a state law to be valid under our Constitution which makes the criminality of an act depend upon the race of the actor.' \textit{McLaughlin v. Florida}, 379 U.S. 184, 198 (concurring opinion). Because I adhere to that belief, I concur in the judgment of the Court.''
\end{quote}
\noindent --- Justice Stewart, concurring, 388 U.S. at 13 (\href{https://supreme.justia.com/cases/federal/us/388/1/}{Justia}; \href{https://www.courtlistener.com/opinion/107480/loving-v-virginia/}{CourtListener}).


\end{document}
